\section{Capture-Avoiding Substitution}
\label{sec:capavoid}

\subsection{Rename, Then Substitute}

The fix is \emph{capture-avoiding substitution}: before substituting into the body of an inner \code{let}, rename its bound variable to a fresh name that appears nowhere else. Renaming a bound variable never changes what a program means (that was the point of the \code{y}/\code{z} example), so we are free to do it whenever it helps.
\begin{lstlisting}[style=repl]
(let x y (let y 2 (+ x y)))
---> (let x y (let g3189 2 (+ x g3189)))    rename the inner y first
---> (let g3189 2 (+ y g3189))              the free y is no longer captured
---> (+ y 2)
---> eval: unbound variable                 same answer as the z version
\end{lstlisting}
The only question is where the fresh name comes from.

\subsection{\texttt{gensym}: Fresh Names}

\code{gensym} (``generate symbol'') is a Racket function that returns a brand-new symbol every time it is called:
\begin{lstlisting}[style=repl]
> (gensym)
'g3189
> (gensym)
'g3190
> (equal? (gensym) (gensym))
#f
\end{lstlisting}
A name chosen by hand, such as \code{'tmp}, does not work: the program might already use \code{tmp}, and then the renaming would create exactly the capture it was meant to avoid.

\begin{definition}[Uninterned symbol]
An ordinary symbol is \emph{interned}: \code{'y} is the same symbol wherever it is written, so \rkt|(equal? 'y 'y)| is \code{\#t}. The symbol returned by \code{gensym} is \emph{uninterned}: it is different from every other symbol, including one spelled the same way.
\end{definition}
Even if a program literally contained the name \code{g3189}, it would not be \code{equal?} to the \code{g3189} that \code{gensym} returned:
\begin{lstlisting}
(define g (gensym))                                ; prints as 'g3189
(equal? g (string->symbol (symbol->string g)))     ; #f
\end{lstlisting}
So a gensym can never be captured by, or capture, any name the user wrote.

\subsubsection{Importing it into \code{plai-typed}} \code{plai-typed} does not provide \code{gensym}, and it cannot check a plain Racket import without knowing its type. \code{typed-in} imports a function from a Racket module and declares its type at the same time:
\begin{lstlisting}
(require (typed-in racket/base [gensym : (-> symbol)]))
\end{lstlisting}
\code{(-> symbol)} means ``takes no arguments and returns a symbol''. The checker trusts this declaration; it does not check Racket's code.

\subsection{\texttt{subst-lazy-v2}}

Only the binding case changes. If the inner \code{let} binds the same name as the one being replaced, stop, as before. Otherwise:
\begin{enumerate}
  \item let \code{z} be \rkt|(gensym)|, a name that appears nowhere else;
  \item rename \code{y} to \code{z} in the body: \rkt|(subst-lazy (idC z) y e2)|. This cannot cause capture, because no \code{let} inside \code{e2} binds the fresh \code{z};
  \item substitute \code{v} for \code{x} in the renamed body.
\end{enumerate}
\begin{lstlisting}
(define (subst-lazy-v2 (v : ArithE) (x : symbol) (e : ArithE)) : ArithE
  (type-case ArithE e
    [numC (n) e]
    [idC (y) (if (equal? x y) v e)]
    [plusC (e1 e2) (plusC (subst-lazy-v2 v x e1) (subst-lazy-v2 v x e2))]
    [timesC (e1 e2) (timesC (subst-lazy-v2 v x e1) (subst-lazy-v2 v x e2))]
    [letC (y e1 e2)
          (if (equal? x y)
              (letC y (subst-lazy-v2 v x e1) e2)
              (let [(z (gensym))]
                (letC z
                      (subst-lazy-v2 v x e1)
                      (subst-lazy-v2 v x (subst-lazy (idC z) y e2)))))]))
\end{lstlisting}
\begin{lstlisting}
(subst-lazy-v2 (idC 'y) 'x (parse '(let y 2 (+ x y))))
;; = (letC 'g3445 (numC 2) (plusC (idC 'y) (idC 'g3445)))
\end{lstlisting}
The inner \code{y} became \code{g3445}, and the \code{y} that came from outside is still free, so evaluation reports it as unbound. Now the \code{y} and \code{z} versions of the program agree.

The lecture's version (\code{lec4a.rkt}) splits the \code{letC} case three ways: \code{x} equals \code{y}; \code{y} is \code{possibly-free?} in \code{v}, so rename; otherwise substitute without renaming. Its \code{possibly-free?} is a stub that always returns \code{\#t}, so it behaves exactly like the two-way code above. Implementing \code{possibly-free?} properly, so that renaming happens only when needed, is left as an exercise (see the free-variable definition in \Cref{sec:letsubst}).

\begin{remark}
Step~2 uses the old \code{subst-lazy}, which may capture in general. It is safe here because the expression being substituted, \rkt|(idC z)|, contains only the fresh name \code{z}, which nothing in \code{e2} can bind.
\end{remark}
